%% LyX 2.4.4 created this file.  For more info, see https://www.lyx.org/.
%% Do not edit unless you really know what you are doing.
\documentclass[11pt,english]{article}
\usepackage[T1]{fontenc}
\usepackage[latin9]{inputenc}
\usepackage{babel}
\usepackage{verbatim}
\usepackage{float}
\usepackage{amsmath}
\usepackage{amssymb}
\usepackage{graphicx}
\usepackage{geometry}
\geometry{verbose,tmargin=1.1in,bmargin=1.1in,lmargin=1.1in,rmargin=1.1in}
\usepackage{setspace}
\usepackage[authoryear,round]{natbib}
\onehalfspacing
\usepackage[bookmarks=false,
 breaklinks=false,pdfborder={0 0 1},backref=section,colorlinks=false]
 {hyperref}

\makeatletter

%%%%%%%%%%%%%%%%%%%%%%%%%%%%%% LyX specific LaTeX commands.
%% Because html converters don't know tabularnewline
\providecommand{\tabularnewline}{\\}

%%%%%%%%%%%%%%%%%%%%%%%%%%%%%% User specified LaTeX commands.
\usepackage{babel}
%
\usepackage{amsfonts}
\newtheorem{prop}{Proposition}
\newtheorem{cor}{Corollary}
\newtheorem{lem}{Lemma}

\setcounter{MaxMatrixCols}{30}
\providecommand{\U}[1]{\protect\rule{.1in}{.1in}}
\newenvironment{proof}[1][Proof]{\noindent\textbf{#1.} }{\ \rule{0.5em}{0.5em}}

\newtheorem{theorem}{Theorem}
\newtheorem{acknowledgement}{Acknowledgement}
\newtheorem{algorithm}[theorem]{Algorithm}
\newtheorem{axiom}[theorem]{Axiom}
\newtheorem{assumption}{Assumption}
\newtheorem{case}[theorem]{Case}
\newtheorem{claim}[theorem]{Claim}
\newtheorem{conclusion}[theorem]{Conclusion}
\newtheorem{condition}[theorem]{Condition}
\newtheorem{conjecture}[theorem]{Conjecture}
\newtheorem{corollary}[theorem]{Corollary}
\newtheorem{criterion}[theorem]{Criterion}
\newtheorem{definition}{Definition}
\newtheorem{example}[theorem]{Example}
\newtheorem{exercise}[theorem]{Exercise}
\newtheorem{lemma}{Lemma}
\newtheorem{notation}[theorem]{Notation}
\newtheorem{problem}[theorem]{Problem}
\newtheorem{proposition}{Proposition}
\newtheorem{remark}[theorem]{Remark}
\newtheorem{solution}[theorem]{Solution}
\newtheorem{summary}[theorem]{Summary}
\usepackage{lscape}
\usepackage{placeins}
\setlength{\emergencystretch}{2.5em}
\sloppy

\makeatother

\begin{document}
\title{Necessity Entrepreneurs,\thanks{We thank}}
\author{David Chivers\thanks{Department of Economics, Durham Business School, DH1 3LB, \protect\protect\href{mailto:david.chivers@durham.ac.uk}{david.chivers@durham.ac.uk}\textbf{ }}
\and Zhigang Feng\thanks{Department of Economics, University of Nebraska, Omaha, NE 68106,
United States, \protect\protect\href{mailto:z.feng2@gmail.com}{z.feng2@gmail.com}.} \and Elisa Keller\thanks{Department of Economics, University of Exeter, \protect\protect\href{mailto:elisa.keller@exeter.ac.uk}{elisa.keller@exeter.ac.uk}}
\and Bo Li\thanks{Department of Economics, Peking University, \protect\protect\href{mailto:lb.harbin@gmail.com\%5C\%2520}{lb.harbin@gmail.com }.} }
\maketitle
\begin{abstract}
TO BE COMPLETED.

\ 

\textbf{JEL Classification}: E23, I10, O40.

\textbf{Keywords}:Occupational Choice, Entrepreneur, Misallocation, 
\end{abstract}
\pagebreak{}

\section{Introduction}

Entrepreneurship is often treated as a driver of innovation, productivity,
and growth. However, not all entrepreneurship has the same economic
role. In many settings, a substantial share of entry is necessity-driven
rather than opportunity-driven. This distinction matters because necessity
entry can reflect weak labor-market outside options rather than comparative
advantage in managing production.

Empirical work typically separates necessity and opportunity entrepreneurship
using survey responses or prior labor-market status. Both approaches
are useful, but each has limits. Self-reported motives are subjective,
while classifications based only on previous employment status can
misclassify agents whose timing is strategic or constrained by insurance
and search frictions. A worker who expects separation may enter early
and look ``opportunity'' in the data; an unemployed household can
enter temporarily while searching and look ``necessity'' even when
underlying talent is high.

This paper provides a model-based definition that maps directly into
counterfactual behavior. We define opportunity entrepreneurs as households
that choose entrepreneurship regardless of whether they enter the
period employed or non-employed. We define necessity entrepreneurs
as households that choose entrepreneurship only when they lack an
employment relationship. The model combines heterogeneous managerial
talent, human-capital heterogeneity, search frictions, and social
insurance, allowing us to quantify how labor-market risk and policy
jointly shape the composition of entrepreneurship.

Our contribution is threefold. First, we provide a structural framework
in which necessity entrepreneurship emerges endogenously from occupational
choice under incomplete insurance. Second, we provide a quantitative
decomposition of how unemployment risk and unemployment insurance
affect entry margins and firm scale across education groups. Third,
we reinterpret necessity entry as a misallocation channel: low outside
options can push households into entrepreneurship even when their
productivity is better used in paid work.

The rest of the paper proceeds as follows. Section 2 discusses related
literature. Section 3 presents the economic environment. Section 4
characterizes household behavior and equilibrium. Section 5 discusses
calibration. Section 6 presents baseline quantitative results. Section
7 studies policy counterfactuals, and Section 8 concludes.

\section{Literature Review}

This paper contributes at the intersection of four literatures: occupational
choice, entrepreneurship under financial frictions, labor-market search,
and heterogeneity in entrepreneurial motives.

The first strand is occupational-choice theory with heterogeneous ability
and outside options. \citet{lucas_1978} provides the canonical
span-of-control benchmark. \citet{evans_jovanovic_1989} and
\citet{hurst_lusardi_2004} show that liquidity constraints
shape who can enter entrepreneurship. \citet{cagetti_denardi_2006}
then quantifies how entrepreneurial frictions and wealth dynamics interact
in general equilibrium.

The second strand studies entrepreneurship and firm dynamics under financial
constraints over the cycle. \citet{buera_2009} shows that credit
frictions affect both entry and post-entry scale. \citet{poschke_2013}
is especially close to our setup: labor-market prospects are central
for occupational choice into entrepreneurship. \citet{sedlacek_sterk_2017}
further show that startup cohorts differ in growth potential over the
business cycle.

The third strand is the search-and-matching labor literature. \citet{mortensen_pissarides_1994}
and \citet{petrongolo_pissarides_2001} provide the benchmark
framework for job-finding and separation dynamics that discipline our
transition structure.

The fourth strand emphasizes entrepreneurial motives and composition.
\citet{fairlie_fossen_2019} distinguish opportunity and necessity
entrepreneurship in the data, while \citet{hurst_pugsley_2011}
show that most small-business activity is not primarily innovation-driven.
\citet{poschke_necessity_snapshot_2013} provides direct evidence
that necessity-based entry is quantitatively important in cross section.
These findings motivate our focus on composition, not only the aggregate
entrepreneurship rate.

Relative to existing work, our contribution is to integrate these channels
in a single heterogeneous-agent framework with explicit state-contingent
occupational choice. This allows us to quantify how unemployment risk
and social insurance reshape the mix of necessity and opportunity entrepreneurship,
firm scale, and aggregate outcomes.

\section{The Model: Economic Environment}

Consider a Lucas (1978) span of control model, where individuals differ
in the level of human capital, $h.$ As in Gomez and Kuhen (2017),
individuals are educated to either the primary $(P)$, secondary $(S)$
or tertiary level $(T)$. The level of education means that a tertiary
educated individuals have a higher draw of human capital, on average,
than secondary educated individuals, who in turn have a higher draw
on average than primary educated individuals. This human capital draw
is positively correlated with managerial ability $x$ as well as labor
productivity $z$.

At the start of the period agents enter as either unemployed $($$i_{w}=0)$,
employed $($$i_{w}=1)$, or an entrepreneur $($$i_{w}=2)$. Individuals
face a choice of whether to become an entrepreneur, to continue working
or to search for work. However, the value of employment is different
depending on whether an agent enters the period as employed $($$i_{w}=1)$
or not. Although there is a risk that an employee may become separated
from a firm with probability $\rho$, this is less risky than searching
for work which has a $p(\theta)$ probability of obtaining employment.
Consequently, the expected value of employment is higher for those
already in employment than those who are not. Additionally, the probability
of separation is lower for individuals with higher levels of human
capital and they are able to command a higher wage due to higher levels
of productivity. So the value of employment, holding other things
equal, is higher for those with higher human capital as well.

The innovation of this model is that the entrepreneurial decision
results in two types of entrepreneurs emerge, opportunity-based entrepreneurs
and necessity-based entrepreneurs. We define opportunity-based entrepreneurs
are those whose value of entrepreneurship is higher, even if they
had employment. Necessity-based entrepreneurs are those who would
not have become entrepreneurs, if they had employment.

\paragraph{Timing: }

The timing of events can be summarized below. Occupational choice
happens at the start of time $t.$ Those choosing to become entrepreneurs
run projects and hire workers on on rolling contracts (i.e. contracts,
will be reassessed each period). Those who were previously employed
within a project can decide to carry on working but carry a risk of
separation with probability $\rho$. If separation occurs these workers
will claim unemployed benefits and be unable to look for work until
the following period. After separation occurs, vacancies are posted
and who were unemployed in the previous period are matched with firms.
Those that are successful and match with with probability $p(\theta)$
go on to earn a wage. Whereas those that are unsuccessful will claim
unemployment benefits and be unable to search again until the following
period.~

Unlike the search and matching literature there is no wage bargaining
between workers and the firm at the end of matching. Instead the wage
pins down the occupational choice decision which determines those
wanting to become entrepreneurs and those who will search for work.
Crucially, the worker only finds out whether they enter production
and earn a wage after their occupational choice decision.

\begin{figure}[H]
\begin{centering}
\caption{Timing for Matched Workers}
\par\end{centering}
\includegraphics[width=0.82\textwidth]{../figures/enter_period_employed_submission} 
\end{figure}

\begin{figure}[H]
\caption{Timing for Unmatched Workers}

\includegraphics[width=0.82\textwidth]{../figures/enter_period_unemployed_submission} 
\end{figure}


\subsection{Preferences, endowments and technology}

\paragraph{Preferences: }

Consumption by an agent in period $t$ is $c_{t}$, with utility given
by $U(c_{t})$.

\paragraph{Endowments: }

Agents are endowed with an initial capital asset, $a_{0}$, which
can be used as an input in production. Each individual is endowed
with human capital shock $h^{i}$. Human capital $h$ is positively
correlated with managerial talent $x$, labor productivity $z$ and
separation rate $\rho.$

\paragraph{Production: }

Firms use labor $n$ and capital $k$ to produce a single consumption
good, $y$. Capital depreciates at a constant rate of $\delta$. Managers
can operate only one project. The functional form of the production
function is: 
\begin{equation}
y=xk^{\alpha}(zn){}^{\gamma}\ \ \mbox{where}\ \alpha,\gamma>0.\ \label{eq:firm_prod}
\end{equation}
Managerial talent $x$ is determined by the Pareto distribution $F(x)=\left[1+\frac{x-\varsigma}{\xi}\right]^{-\vartheta_{i}}$
where $\vartheta$ is the shape parameter, $\varsigma$ the location
parameter and $\xi$ is the scale parameter.

\paragraph{Factor remuneration: }

Firms rent capital at the common market rate $r$.

\paragraph{Government:}

The government runs a balanced budget each period and provides unemployment
insurance which is financed through lump sum taxation $\tau$, which
is shared by the employee $(1-\psi)\tau$ and employer $\psi\tau$.
This program guarantees each household an unemployment benefit of
$b$.

\subsection{The Firm's problem}

The firm's problem is: 
\begin{equation}
\max_{n,k}Xk^{\alpha}n^{\gamma}-(1+\tau\psi)wzn-rk-\kappa\frac{\rho zn}{q(\theta_{t})}\label{eq:max_firm}
\end{equation}
employees must pay a share of unemployment insurance $\tau$ to all
workers who produce output. Firms have to pay search costs $\kappa$
for all vacancies posted. The number of vacancies posted is designed
to replace any lost workers, $\rho zn.$ This, however, is affected
by the firm-hiring rate $q(\theta_{t})=\bar{m}\left(\frac{v_{t}}{u_{t}}\right)^{-\mu}$,
where aggregate vacancies $v_{t}$ is simply the aggregate level of
all those who have separate from firms $\sum\rho zn$ within the period
and $u_{t}$ is the total unemployed looking for work within that
period (see appendix).

When the market is tight and the hiring rate is low, firms have to
search more and thus face higher search costs in order to replace
lost workers. $\text{This is because }0<q(\theta_{t})<1,$so a lower
$q(\theta_{t})$ increases $\frac{\rho zn}{q(\theta_{t})}$. If the
firm lose 10 workers, they want to hire 10 workers to replace them.
But if the hiring rate, $q(\theta_{t}),$ is 50\% then they will face
double the amount of search costs than if $q(\theta_{t})$ were 100\%.
The upshot of this is that when markets are tight, the size of the
desired firm size $n$ decreases ex ante , as the cost of replacing
workers is high and so entrepreneurs would prefer to run smaller firms.

To avoid issues of directed search, we assume that employees are randomly
matched with firms (as in Chivers et al., 2015) and each firm employees
a distribution of workers across the human capital spectrum such that
each firm employs the market average.

\section{Optimal behavior and equilibrium}

In any time period $t$ agents are distinguished by $(h,a,x,i,z)$.
The timing of the economy is given as follows. 
\begin{enumerate}
\item Households enter each new period with assets $a$ and an employment
status $i_{w}$ 
\item Idiosyncratic shocks $h$ is drawn which determines $x$ (managerial
ability) and $z$ (labor productivity), $\rho$ seperation rate conditional
on education level, $P,S$ or $T.$ 
\item Households make an occupation decision: entrepreneur ($\mathcal{I}_{e}=1$)
or worker ($\mathcal{I}_{e}=0$). 
\item Capital and labor markets clear. A fraction $\rho$ of employees separate
from firms and~$p(\theta)$ are hired to replace them 
\item Households choose: consumption ($c$), borrowing/saving ($a'$) 
\end{enumerate}

\subsection{Firm manager}

Firms are distinguished by their productivity realization $x$. Agents
with sufficient ability to become managers choose the level of capital
and the number of employees to maximize profit subject to a technological
constraint. In order to simplify the exposition, first consider the
problem of a manager with talent $x^{i}$ for a given level of capital
$k$ (i.e., labor input choice only): 
\begin{equation}
\max_{zn}xk^{\alpha}(zn){}^{\gamma}-zw(1+\psi\tau)n-rk-\kappa\frac{\rho zn}{q(\theta)}\label{maxP}
\end{equation}

The first order conditions are:

\begin{equation}
\gamma xk^{\alpha}(zn)^{\gamma-1}=w(1+\psi\tau)+\frac{\rho\kappa}{q(\theta)}\label{maxP-1}
\end{equation}

\begin{equation}
n^{*}(k,x,\widetilde{w})=\left[\frac{\gamma xk^{\alpha}}{w(1+\psi\tau)+\frac{\rho\kappa}{q(\theta)}}\right]^{\frac{1}{1-\gamma}}\label{FOC}
\end{equation}

Substituting (\ref{FOC}) into (\ref{maxP}) yields the manager's
profit function for a given level of capital: 
\begin{equation}
y(k,x,\widetilde{w})=xk^{\alpha}(zn^{*})^{\gamma}-zn^{*}w(1+\psi\tau)-\kappa\frac{\rho zn^{*}}{q(\theta)}\label{profitj}
\end{equation}


\subsubsection{Capital}

Initially we assume capital markets function perfectly.

\subsection{Entrepreneurs}

Entrepreneurs maximize expected discounted utility of consumption
\[
\max_{\{c_{t},a_{t+1}\}}\mathbb{E}\sum_{t=0}^{\infty}\beta^{t}U(c_{t})
\]
subject to the following budget constraint: 
\begin{align}
c+a' & \leq(1+r-\delta)a+\nu-\tau_{\pi}\label{eq:c_worker-3}
\end{align}

\[
a'\geq-\bar{a}
\]


\subsection{Employed}

The workforce (previous periods employed who choose to reapply for
work) maximize expected discounted utility of consumption 
\[
\max_{\{c_{t},a_{t+1}\}}\mathbb{E}\sum_{t=0}^{\infty}\beta^{t}U(c_{t})
\]
subject to the following budget constraint: 
\begin{align}
c+a' & \leq(1+r-\delta)a+(1-\rho)\left(wz(1-(1-\psi)\tau_{y})\right)+\rho b\label{eq:c_worker}
\end{align}

\[
a'\geq-\bar{a}
\]
where $b$ is the unemployment benefit.

\subsection{Unemployed}

The unemployed maximize expected discounted utility of consumption
\[
\max_{\{c_{t},a_{t+1}\}}\mathbb{E}\sum_{t=0}^{\infty}\beta^{t}U(c_{t})
\]
subject to the following budget constraint: 
\begin{align}
c+a' & \leq(1+r-\delta)a+p(\theta)\left(wz(1-(1-\psi)\tau_{y})\right)+(1-p(\theta))b\label{eq:c_worker-2}
\end{align}

\[
a'\geq-\bar{a}
\]
where $b$ is the unemployment benefit.

\subsection{Government }

The government runs a balanced budget with a lump-sum taxes $T(y)=(y-\lambda_{p}y^{1-\tau_{p}})$

\subsection{The household's problem}

Let $i_{w}$ denote the household's labor market status in the previous
period, with $i_{w}=2$ indicating being entrepreneur, $i_{w}=1$
being employed and $i_{w}=0$ being unemployed. Let $\mathcal{I}_{e}$
indicate occupational choice, where $\mathcal{I}_{e}=1$ if the household
is an entrepreneur and $\mathcal{I}_{e}=0$ if the household is a
worker.

Let denote $\mathbb{E}_{\rho}$ the expected value considering the
possibility of being separated from the job as employed; $\mathbb{E}_{p(\theta)}$
the expected value considering the possibility of finding new job
if previously unmatched with a firm (unemployed former entrepreneurs).
$V_{w}$ is the contemporary utility as a worker, which varies with
the employment status. $\iota_{i_{w}=1}$ is an indicator function.
We can write the household's problem recursively as follows. More
specifically, 
\begin{align*}
\mathbb{E}_{p(\theta)}\left(V_{w}+\beta\mathbb{E}\mathbf{V}(a',x',i'_{w})\right) & =p(\theta)\left(V_{n}+\beta\mathbb{E}\mathbf{V}(a',x',i'_{w}=1)\right)\\
 & +(1-p(\theta))\left(V_{u}+\beta\mathbb{E}\mathbf{V}(a',x',i'_{w}=0)\right)
\end{align*}

\begin{align*}
\mathbb{E}_{\rho}\left(V_{w}+\beta\mathbb{E}\mathbf{V}(a',x',i'_{w})\right) & =\rho\left(V_{u}+\beta\mathbb{E}\mathbf{V}(a',x',i'_{w}=0)\right)\\
 & +(1-\rho)\left(V_{n}+\beta\mathbb{E}\mathbf{V}(a',x',i'_{w}=1)\right)
\end{align*}
Notice that for agents who recently separated from the firm cannot
rematch with a new firm in the current period.

\begin{align*}
\mathbf{V}(a,h,i_{w}) & =\max_{\left\{ a',c,\mathcal{I}_{e}\right\} }\left\{ \begin{array}{c}
\iota_{i_{w}=0}\left[(1-\mathcal{I}_{e})\mathbb{E}_{p(\theta)}\left(V_{w}+\beta\mathbb{E}\mathbf{V}(a',x',i'_{w})\right)+\mathcal{I}_{e}\left(V_{e}+\beta\mathbb{E}\mathbf{V}(a',x',i'_{w}=0)\right)\right]+\\
\iota_{i_{w}=1}\left[(1-\mathcal{I}_{e})\mathbb{E}_{\rho}\left(V_{w}+\beta\mathbb{E}\mathbf{V}(a',x',i'_{w})\right)+\mathcal{I}_{e}\left(V_{e}+\beta\mathbb{E}\mathbf{V}(a',x',i'_{w}=0)\right)\right]+\\
\iota_{i_{w}=2}\left[(1-\mathcal{I}_{e})\mathbb{E}_{p(\theta)}\left(V_{w}+\beta\mathbb{E}\mathbf{V}(a',x',i'_{w})\right)+\mathcal{I}_{e}\left(V_{e}+\beta\mathbb{E}\mathbf{V}(a',x',i'_{w}=0)\right)\right]
\end{array}\right\} 
\end{align*}
subject to

\begin{align}
c+a' & \leq(1-\tilde{{r}}-\delta)a+inc\label{eq:c_worker-1}
\end{align}
\begin{align*}
inc & =\begin{cases}
ra+\widetilde{w}_{n} & \mbox{ if \ensuremath{\mathcal{I}_{e}=0,} \ensuremath{i_{w}=1}}\\
ra+\widetilde{w}_{u} & \mbox{ if \ensuremath{\mathcal{I}_{e}=0,} \ensuremath{i_{w}\neq1}}\\
ra+\nu(k,x;\tilde{r},w,\tau) & \mbox{ if \ensuremath{\mathcal{I}_{e}=1}}
\end{cases}
\end{align*}

where $inc$ is the earnings of the worker or entrepreneur, $\widetilde{w}_{n}=(1-\rho)\left\{ zw(1-(1-\psi)\tau)\right\} +\rho b$
and $\widetilde{w}_{u}=p(\theta)\left\{ zw(1-(1-\psi)\tau)\right\} +\{1-p(\theta)\}b$.

\subsection{Equilibrium}

Let $\hat{h}$ be the critical value of human capital that determines
occupational choice which is affected by both employment status and
education level. If $h\geq\hat{h}$ then individuals choose entrepreneurship
over becoming a worker. The critical value of human capital is always
higher for those who have employment than those who are unmatched
with a firm $\hat{h}_{i_{w}\neq1}<\hat{h}_{i_{w}=1}$. We define necessity
entrepreneurship as those who would not have become an entrepreneur
\textit{if} they had an employment contract: 
\[
\hat{h}_{i_{w}\neq1}<h<\hat{h}_{i_{w}=1}.
\]
We define opportunity-driven entrepreneurship as those who would become
an entrepreneur \textit{regardless} of circumstance: 
\[
h\geq\hat{h}_{i_{w}=1}\geq\hat{h}_{i_{w}\neq1}
\]
As there are three education levels, there are a total of six critical
values of human capital: two for primary educated matched or unmatched,
$\hat{h}_{i_{w}=1}^{P}\text{ and }\hat{h}_{i_{w}\neq1}^{P}$ two for
secondary educated, $\hat{h}_{i_{w}=1}^{S}\text{ and }\hat{h}_{i_{w}\neq1}^{S},$
and two for tertiary educated, $\hat{h}_{i_{w}=1}^{T}\text{ and }\hat{h}_{i_{w}\neq1}^{T}$.

\section{Calibration}
\begin{center}
\begin{table}[htbp]
\caption{Parameter values, baseline economy}

\centering{}\resizebox{\textwidth}{!}{{\small{}%
\begin{tabular}{cccc}
\hline 
{\small Parameters } & {\small Values } & {\small Description } & {\small Comments/observations/target}\tabularnewline
\hline 
{\small$\beta$ } & {\small$0.94$ } & {\small Discount factor } & \tabularnewline
{\small$\sigma$ } & {\small$2.0$ } & {\small Risk aversion } & {\small$u(c)=\frac{c^{1-\sigma}}{1-\sigma}$}\tabularnewline
\hline 
{\small$\delta$ } & {\small$0.06$ } & {\small Capital depreciation } & \tabularnewline
{\small$\alpha$ } & {\small$0.3207$ } & {\small Capital share } & \tabularnewline
{\small$\gamma$ } & {\small$0.4693$ } & {\small Labour Share } & {\small Frac. of entrepreneurs}\tabularnewline
\hline 
{\small$h$ } & {\small$\{0.109,0.581,0.310\}$ } & {\small Share of population by $h$ } & {\small Barro and Lee (2013)}\tabularnewline
{\small$\vartheta_{e}$ } & {\small$\{0.967,1.007,1.156\}$ } & {\small Shape parameter by $h$ } & {\small Mean firm size, Gomez and Kuhen (2017)}\tabularnewline
{\small$\varsigma$ } & {\small$2.02655$ } & {\small Location parameter } & {\small$F(x)=\left[1+\frac{x-\varsigma}{\xi}\right]^{-\alpha_{i}}$}\tabularnewline
{\small$\xi$ } & {\small$0.1115$ } & {\small Scale parameter } & {\small Gomez and Kuhen (2017)}\tabularnewline
{\small$z$ } & {\small$\{1.0,1.31,1.61\}$ } & {\small Wage premium by $h$ } & {\small Goldin and Katz (2009)}\tabularnewline
{\small$\rho$ } & {\small$\{0.0383,0.0288,0.0136\}$ } & {\small Job separation rate by $h$ } & {\small Cairo and Cajner (2016)}\tabularnewline
\hline 
{\small$\bar{m}$ } & {\small$0.8$ } & {\small Matching efficiency } & {\small Lubik, 2013 $p(\theta)=\bar{m}\left(\frac{v}{u}\right)^{1-\mu}$}\tabularnewline
{\small$\mu$ } & {\small$[0.5,0.7]$ } & {\small Matching elasticity } & {\small Petrongolo and Pissarides (2001)}\tabularnewline
{\small$\kappa$ } & {\small$[0.03,0.045]\cdot w$ } & {\small Search costs } & {\small (Hagedorn and Manovskii, 2008)}\tabularnewline
\hline 
{\small$\tau_{\pi}$ } & {\small$0$ } & {\small Corporation tax } & {\small Set as zero for now}\tabularnewline
{\small$\tau_{y}$ } & {\small$0.151$ } & {\small Income Tax } & {\small$T(y)=(y-\lambda_{p}y^{1-\tau_{p}})$}\tabularnewline
{\small$\lambda_{p}$ } &  & {\small Income Tax } & {\small Determined in equm.}\tabularnewline
{\small$b$ } & {\small$0.40\cdot w$, } & {\small Social security benefits } & {\small Shimmer and Krause and Lubik (2005)}\tabularnewline
{\small$\psi$ } & {\small$0$ } & {\small Share } & \tabularnewline
{\small$\Delta$ } & {\small$0$ } & {\small Credit spread } & \tabularnewline
{\small$\bar{a}$ } & {\small$0$ } & {\small Borrowing constraint } & \tabularnewline
 &  &  & \tabularnewline
\hline 
\end{tabular}}} 
\end{table}
\par\end{center}

\section{Baseline Results}

The baseline equilibrium shows moderate entrepreneurship (0.1472 overall),
with substantial heterogeneity by education. Entry is highest for
high education (0.1618) and lowest for middle education (0.1398),
while average labor demand is 21.9129 workers and rises steeply with
skill from 9.7237 (low) to 33.5470 (high). This pattern indicates
that high-education entrepreneurs operate at much larger scale even
when their entry rate is only modestly above other groups. Labor-market
transitions also suggest strong persistence in employment (Employed
$\rightarrow$ Work = 0.740) and a small unemployment-to-entrepreneurship
channel (0.002), so necessity entry is present but limited in the
baseline.

\begin{table}[H]
\caption{Baseline equilibrium outcomes}
\label{tab:baseline}
\centering
\small
\begin{tabular}{|l|c|}
\hline
 & Baseline \\
\hline
UI & 0.40 \\
\hline
Interest rate $r$ & 0.0989 \\
\hline
Wage $w$ & 2.3870 \\
\hline
Job-finding rate $\theta$ & 0.5063 \\
\hline
Income tax $\tau_y$ & 0.0377 \\
\hline
Convergence metric & 0.0010 \\
\hline
Govt consumption floor pay & 332.3780 \\
\hline
Entrepreneurship (all edu) & 0.1472 \\
\hline
Entrepreneurship (low edu) & 0.1457 \\
\hline
Entrepreneurship (med edu) & 0.1398 \\
\hline
Entrepreneurship (high edu) & 0.1618 \\
\hline
Entr $\rightarrow$ Entr & 0.145 \\
\hline
Entr $\rightarrow$ Work & 0.000 \\
\hline
Emp $\rightarrow$ Entr & 0.000 \\
\hline
Emp $\rightarrow$ Work & 0.740 \\
\hline
Unemp $\rightarrow$ Entr & 0.002 \\
\hline
Unemp $\rightarrow$ Work & 0.113 \\
\hline
Avg labor demand $n$ & 21.9129 \\
\hline
Labor demand $n$ (low) & 9.7237 \\
\hline
Labor demand $n$ (med) & 17.9923 \\
\hline
Labor demand $n$ (high) & 33.5470 \\
\hline
Std. dev. labor demand & 27.9509 \\
\hline
Per-capita output & 4.5316 \\
\hline
Output per entrepreneur & 31.3760 \\
\hline
\end{tabular}
\end{table}


\section{Quantitative Analysis}

This section evaluates each policy experiment relative to the baseline
and reports how unemployment insurance design and tax closure affect
entry, scale, and output.

\FloatBarrier

\subsection{Low unemployment insurance with endogenous tax}

This subsection reports the low-UI endogenous-tax experiment against
the baseline.

\begin{table}[H]
\caption{Policy experiment: Baseline vs low unemployment insurance with endogenous tax}
\label{tab:exp_lowui_endogtax}
\centering
\small
\begin{tabular}{|l|c|c|}
\hline
 & Baseline & Low UI (Endogenous Tax) \\
\hline
UI & 0.40 & 0.20 \\
\hline
Interest rate $r$ & 0.0989 & 0.0982 \\
\hline
Wage $w$ & 2.3870 & 2.4036 \\
\hline
Job-finding rate $\theta$ & 0.5063 & 0.5137 \\
\hline
Income tax $\tau_y$ & 0.0377 & 0.0182 \\
\hline
Convergence metric & 0.0010 & 0.0009 \\
\hline
Govt consumption floor pay & 332.3780 & 218.5260 \\
\hline
Entrepreneurship (all edu) & 0.1472 & 0.1542 \\
\hline
Entrepreneurship (low edu) & 0.1457 & 0.1542 \\
\hline
Entrepreneurship (med edu) & 0.1398 & 0.1472 \\
\hline
Entrepreneurship (high edu) & 0.1618 & 0.1672 \\
\hline
Entr $\rightarrow$ Entr & 0.145 & 0.151 \\
\hline
Entr $\rightarrow$ Work & 0.000 & 0.000 \\
\hline
Emp $\rightarrow$ Entr & 0.000 & 0.000 \\
\hline
Emp $\rightarrow$ Work & 0.740 & 0.738 \\
\hline
Unemp $\rightarrow$ Entr & 0.002 & 0.003 \\
\hline
Unemp $\rightarrow$ Work & 0.113 & 0.107 \\
\hline
Avg labor demand $n$ & 21.9129 & 20.6400 \\
\hline
Labor demand $n$ (low) & 9.7237 & 9.0788 \\
\hline
Labor demand $n$ (med) & 17.9923 & 16.8165 \\
\hline
Labor demand $n$ (high) & 33.5470 & 31.8708 \\
\hline
Std. dev. labor demand & 27.9509 & 27.1502 \\
\hline
Per-capita output & 4.5316 & 4.8523 \\
\hline
\% vs baseline (pc output) & 0.00\% & +7.08\% \\
\hline
Output per entrepreneur & 31.3760 & 31.1302 \\
\hline
\% vs baseline (productivity) & 0.00\% & -0.78\% \\
\hline
\end{tabular}
\end{table}


Lower UI increases entrepreneurship and per-capita output relative
to baseline, while output per entrepreneur is slightly lower.

\FloatBarrier

\subsection{No unemployment insurance with endogenous tax}

This subsection reports the no-UI endogenous-tax experiment against
the baseline.

\begin{table}[H]
\caption{Policy experiment: Baseline vs no unemployment insurance with endogenous tax}
\label{tab:exp_noui_endogtax}
\centering
\small
\begin{tabular}{|l|c|c|}
\hline
 & Baseline & No UI (Endogenous Tax) \\
\hline
UI & 0.40 & 0.05 \\
\hline
Interest rate $r$ & 0.0989 & 0.0969 \\
\hline
Wage $w$ & 2.3870 & 2.6273 \\
\hline
Job-finding rate $\theta$ & 0.5063 & 0.4965 \\
\hline
Income tax $\tau_y$ & 0.0377 & 0.0008 \\
\hline
Convergence metric & 0.0010 & 0.0003 \\
\hline
Govt consumption floor pay & 332.3780 & 89.6635 \\
\hline
Entrepreneurship (all edu) & 0.1472 & 0.2163 \\
\hline
Entrepreneurship (low edu) & 0.1457 & 0.2481 \\
\hline
Entrepreneurship (med edu) & 0.1398 & 0.2180 \\
\hline
Entrepreneurship (high edu) & 0.1618 & 0.2018 \\
\hline
Entr $\rightarrow$ Entr & 0.145 & 0.212 \\
\hline
Entr $\rightarrow$ Work & 0.000 & 0.000 \\
\hline
Emp $\rightarrow$ Entr & 0.000 & 0.000 \\
\hline
Emp $\rightarrow$ Work & 0.740 & 0.676 \\
\hline
Unemp $\rightarrow$ Entr & 0.002 & 0.004 \\
\hline
Unemp $\rightarrow$ Work & 0.113 & 0.108 \\
\hline
Avg labor demand $n$ & 21.9129 & 11.5534 \\
\hline
Labor demand $n$ (low) & 9.7237 & 4.5716 \\
\hline
Labor demand $n$ (med) & 17.9923 & 8.7990 \\
\hline
Labor demand $n$ (high) & 33.5470 & 19.1706 \\
\hline
Std. dev. labor demand & 27.9509 & 18.5538 \\
\hline
Per-capita output & 4.5316 & 4.6991 \\
\hline
\% vs baseline (pc output) & 0.00\% & +3.70\% \\
\hline
Output per entrepreneur & 31.3760 & 25.1883 \\
\hline
\% vs baseline (productivity) & 0.00\% & -19.72\% \\
\hline
\end{tabular}
\end{table}


Removing UI under endogenous tax strongly raises entry, but average
labor demand and output per entrepreneur decline materially.

\FloatBarrier

\subsection{No unemployment insurance with fixed tax}

This subsection reports the no-UI fixed-tax experiment against the
baseline.

\begin{table}[H]
\caption{Policy experiment: Baseline vs no unemployment insurance with fixed tax}
\label{tab:exp_noui_fixedtax}
\centering
\small
\begin{tabular}{|l|c|c|}
\hline
 & Baseline & No UI (Fixed Tax) \\
\hline
UI & 0.40 & 0.10 \\
\hline
Interest rate $r$ & 0.0989 & 0.0977 \\
\hline
Wage $w$ & 2.3870 & 2.4061 \\
\hline
Job-finding rate $\theta$ & 0.5063 & 0.5099 \\
\hline
Income tax $\tau_y$ & 0.0377 & 0.0386 \\
\hline
Convergence metric & 0.0010 & 0.0012 \\
\hline
Govt consumption floor pay & 332.3780 & 192.3750 \\
\hline
Entrepreneurship (all edu) & 0.1472 & 0.1530 \\
\hline
Entrepreneurship (low edu) & 0.1457 & 0.1507 \\
\hline
Entrepreneurship (med edu) & 0.1398 & 0.1455 \\
\hline
Entrepreneurship (high edu) & 0.1618 & 0.1680 \\
\hline
Entr $\rightarrow$ Entr & 0.145 & 0.151 \\
\hline
Entr $\rightarrow$ Work & 0.000 & 0.000 \\
\hline
Emp $\rightarrow$ Entr & 0.000 & 0.000 \\
\hline
Emp $\rightarrow$ Work & 0.740 & 0.737 \\
\hline
Unemp $\rightarrow$ Entr & 0.002 & 0.002 \\
\hline
Unemp $\rightarrow$ Work & 0.113 & 0.110 \\
\hline
Avg labor demand $n$ & 21.9129 & 20.7240 \\
\hline
Labor demand $n$ (low) & 9.7237 & 9.2638 \\
\hline
Labor demand $n$ (med) & 17.9923 & 16.9940 \\
\hline
Labor demand $n$ (high) & 33.5470 & 31.7442 \\
\hline
Std. dev. labor demand & 27.9509 & 27.1523 \\
\hline
Per-capita output & 4.5316 & 4.5261 \\
\hline
\% vs baseline (pc output) & 0.00\% & -0.12\% \\
\hline
Output per entrepreneur & 31.3760 & 27.7709 \\
\hline
\% vs baseline (productivity) & 0.00\% & -11.49\% \\
\hline
\end{tabular}
\end{table}


With fixed tax and no UI, entry rises modestly, per-capita output
is roughly unchanged, and productivity per entrepreneur declines.

\FloatBarrier

\section{Conclusion}

This paper develops a high-level quantitative framework for studying
necessity entrepreneurship as an equilibrium outcome of occupational
choice under labor market frictions. By defining necessity and opportunity
entrepreneurship through counterfactual choices, the model separates
compositional shifts in entry from changes in aggregate entrepreneurship
rates.

The main message is that unemployment risk and social insurance jointly
shape who becomes an entrepreneur. Reducing insurance can raise entrepreneurship
on paper, but part of this increase reflects necessity entry associated
with lower average firm scale. Conversely, stronger insurance can
reduce necessity entry and improve sorting, even when effects on total
entrepreneurship are modest. These results support evaluating entrepreneurship
policy through composition, not levels alone.

Several extensions remain important for the next draft. First, the
credit friction wedge and job-market experiments should be fully completed
and integrated into a unified welfare comparison. Second, calibration
choices that are still open should be aligned with the code implementation
and reported transparently. Third, the literature discussion can be
broadened further once the project-specific verified PDF library is
restored.

\section*{\newpage}

\section*{Appendix}

\section*{Computation}

Given the values for parameters, the discretized spaces $\Omega_{a}$
for $a$, $\Omega_{z}$ for $z$ and $\Omega_{x}$ for $x$. Each
agent draw a human capital level $h$ (interpreted as educational
attainment) at the beginning of the economy, which will be fixed henceforth.
$h$ will determine their labor productivity $z\in\Omega_{z}$, the
separation rate $\rho$, and managerial ability $x\in\Omega_{x}$.
Let $i_{w}=0$ indicating being entrepreneur, $i_{w}=1$ being employed
and $i_{w}=2$ being unemployed. The numerical algorithm works as
follows. 
\begin{enumerate}
\item Set a tolerance $\epsilon>0$. 
\item Guess $\Phi^{0}=\left(r^{0},w^{0},\rho^{0},\tau_{y}^{0}\right)$.
Solve for optimal household behavior: 
\[
f:\left(\varphi;\Phi\right)\rightarrow\left(c,a',\mathcal{I}_{e},n,k\right),
\]
where $\varphi=\left\{ h,z,x,a,i_{w}\right\} $. We will use the method
of value function iteration as follows.
\begin{enumerate}
\item Guess value function $\mathbf{V}{}^{0}\left(\varphi;\Phi^{0}\right)$
and policy functions $f^{0}\left(\varphi;\Phi^{0}\right)$. 
\item Calculate the cutoff values of $x_{u}^{*}$ (for previously unemployed
worker, or entrepreneur) and $x_{n}^{*}$ (for previously employed
worker) via the following conditions. 
\begin{align}
\widetilde{w}_{u}= & \pi(x_{u}^{*})=\max_{n}\left\{ \left(x_{u}^{*}\right)^{1-\alpha-\gamma}k^{\alpha}n{}^{\gamma}-w(1+\psi\tau)n-rk-\frac{\kappa zn}{q(\theta)}\right\} \label{eq:xstar_u}\\
\widetilde{w}_{n}= & \pi(x_{n}^{*})=\max_{n}\left\{ \left(x_{n}^{*}\right)^{1-\alpha-\gamma}k^{\alpha}n{}^{\gamma}-w(1+\psi\tau)n-rk-\frac{\kappa zn}{q(\theta)}\right\} \label{eq:xstar_n}
\end{align}
where $\widetilde{w}_{n}$ and $\widetilde{w}_{u}$ denote the expected
compensation for the employed worker and unemployed worker respectively.
More specifically, 
\begin{align}
\widetilde{w}_{n} & =(1-\rho)\left\{ zw(1-(1-\psi)\tau)\right\} +\rho b\label{eq:income_n}\\
\widetilde{w}_{u} & =p(\theta)\left\{ zw(1-(1-\psi)\tau)\right\} p(\theta))b.\label{eq:income_u}
\end{align}
Noting that the separation rate $\rho$ is pre-determined by the individual's
initial human capital level $h$. 
\item Update value and policy functions: 
\begin{align}
\mathbf{V}^{1}(h,z,x,a,i_{w})=\label{eq:V_overall}\\
\max_{\left\{ a',c,\mathcal{I}_{e}\right\} }\left\{ \begin{array}{c}
\iota_{i_{w}=0}\cdot\left[\begin{array}{c}
(1-\mathcal{I}_{e})\cdot\mathbb{E}_{p(\theta)}\left(V_{w}+\beta\mathbb{E}\mathbf{V}^{0}(h,z',x',a',i'_{w})\right)\\
+\mathcal{I}_{e}\cdot\left(V_{e}+\beta\mathbb{E}\mathbf{V}^{0}(h,z',x',a',i'_{w}=2)\right)
\end{array}\right]+\\
\iota_{i_{w}=1}\cdot\left[\begin{array}{c}
(1-\mathcal{I}_{e})\cdot\mathbb{E}_{\rho}\left(V_{w}+\beta\mathbb{E}\mathbf{V}^{0}(h,z',x',a',i'_{w})\right)\\
+\mathcal{I}_{e}\cdot\left(V_{e}+\beta\mathbb{E}\mathbf{V}^{0}(h,z',x',a',i'_{w}=2)\right)
\end{array}\right]+\\
\iota_{i_{w}=2}\cdot\left[\begin{array}{c}
(1-\mathcal{I}_{e})\cdot\mathbb{E}_{p(\theta)}\left(V_{w}+\beta\mathbb{E}\mathbf{V}^{0}(h,z',x',a',i'_{w})\right)\\
+\mathcal{I}_{e}\cdot\left(V_{e}+\beta\mathbb{E}\mathbf{V}^{0}(h,z',x',a',i'_{w}=2)\right)
\end{array}\right]
\end{array}\right\} \nonumber \\
f^{1}\left(\varphi;\Phi^{0}\right)=\arg\max\mathbf{V}^{1}\left(\varphi;\Phi^{0}\right)\nonumber 
\end{align}
where $\mathbb{E}_{p(\theta)}\left(V_{w}+\beta\mathbb{E}\mathbf{V}(a',x',i'_{w})\right)$
denotes the expected payoff for the individual who is previously unemployed
worker, or entrepreneur, and decides to look for paid job, while $\mathbb{E}_{\rho}\left(V_{w}+\beta\mathbb{E}\mathbf{V}(a',x',i'_{w})\right)$
represents the expected payoff for the previously employed worker
who decide to stay in workplace. 
\begin{align}
\mathbb{E}_{p(\theta)}\left(V_{w}+\beta\mathbb{E}\mathbf{V}(h,z',x',a',i'_{w})\right) & =p(\theta)\left(V_{n}+\beta\mathbb{E}\mathbf{V}(h,z',x',a',i'_{w}=1)\right)\label{eq:V_u}\\
 & +(1-p(\theta))\left(V_{u}+\beta\mathbb{E}\mathbf{V}(h,z',x',a',i'_{w}=0)\right)\nonumber \\
\mathbb{E}_{\rho}\left(V_{w}+\beta\mathbb{E}\mathbf{V}(h,z',x',a',i'_{w})\right) & =(1-\rho)\left(V_{n}+\beta\mathbb{E}\mathbf{V}(h,z',x',a',i'_{w}=1)\right)\label{eq:V_w}\\
 & +\rho\left(V_{u}+\beta\mathbb{E}\mathbf{V}(h,z',x',a',i'_{w}=0)\right)\nonumber 
\end{align}
To simplify the analysis, we assume that the individual who was previously
employed and newly separated from the firm is not allowed to search
for job at the current period. The measure for the individuals who
are looking for jobs, $u_{t}$, the total vacancies $v_{t}$, and
job matching rate $\theta_{n,t}$ are given as follows. 
\begin{align}
u_{t} & =\left(\int_{x}\iota_{i_{w}=0}+\int_{x}\iota_{i_{w}=2}\right)\cdot\int_{0}^{x_{u}^{*}}\label{eq:measure_u}\\
v_{t} & =\int_{x}\iota_{i_{w}=1}\cdot\rho\label{eq:measure_v}\\
\theta_{n,t} & =\bar{m}\left(\frac{v_{t}}{u_{t}}\right)^{1-\mu}\label{eq:measure_matching}
\end{align}
Note that we can write the above value functions explicitly as a system
of sub-problems: $P_{x<x_{u}^{*}}^{i_{w}=0}$ , $P_{x\geq x_{u}^{*}}^{i_{w}=0}$,
$P_{x<x_{n}^{*}}^{i_{w}=1}$ , $P_{x\geq x_{n}^{*}}^{i_{w}=1}$, $P_{x<x_{u}^{*}}^{i_{w}=2}$
, $P_{x\geq x_{u}^{*}}^{i_{w}=2}$, which correspond to the dynamic
programing problem for i). previously self-employed with $x<x_{u}^{*}$;
ii). previously self-employed with $x\geq x_{u}^{*}$; iii). previously
employed with $x<x_{n}^{*}$; iv). previously employed with $x\geq x_{n}^{*}$;
v). previously unemployed with $x<x_{u}^{*}$; vi). previously unemployed
with $h\geq h_{u}^{*}$. To economize the space, we only provide the
details for $P_{x<x_{u}^{*}}^{i_{w}=0}$ , and $P_{x\geq x_{u}^{*}}^{i_{w}=0}$
as follows. \\
 $\left(P_{x<x_{u}^{*}}^{i_{w}=0}\right)$: 
\begin{equation}
\mathbf{V}^{1}(h,z,x<x_{u}^{*},a,i_{w}=0)=\max_{a'_{n},a'_{u}}\left\{ \begin{array}{c}
p(\theta)\left(u(c_{n})+\beta\mathbb{E}\mathbf{V}^{0}(h,z',x',a'_{n},i'_{w}=1)\right)\\
+(1-p(\theta))\left(u(c_{u})+\beta\mathbb{E}\mathbf{V}^{0}(h,z',x',a'_{u},i'_{w}=0)\right)
\end{array}\right\} \label{eq:V_iw0_xl}
\end{equation}
s.t. 
\begin{align}
c_{n}+a'_{n} & \leq(1-\tilde{r}-\delta)a+zw(1-(1-\psi)\tau),\text{if matched with a firm}\label{eq:bdg_iw0_xl_1}\\
c_{u}+a'_{u} & \leq(1-\tilde{r}-\delta)a+b,\text{if unemployed}\label{eq:bdg_iw0_xl_2}
\end{align}
$\left(P_{x\geq x_{u}^{*}}^{i_{w}=0}\right)$: 
\begin{equation}
\mathbf{V}^{1}(h,z,x\geq x_{u}^{*},a,i_{w}=0)=\max_{a_{e}'}\left\{ u(c_{e})+\beta\mathbb{E}\mathbf{V}^{0}(h,z',x',a_{e}',i'_{w}=w)\right\} \label{eq:V_iw0_xh}
\end{equation}
s.t. 
\begin{equation}
c_{e}+a_{e}'\leq(1-\tilde{r}-\delta)a+\nu(k,x;\tilde{r},w,\tau)\label{eq:bdg_iw0_xh}
\end{equation}
where 
\begin{equation}
\nu(k,x;\tilde{r},w,\tau)=\max_{n}\left\{ x^{1-\alpha-\gamma}k^{\alpha}n{}^{\gamma}-w(1+\psi\tau)n-rk-\frac{\kappa\rho n}{q(\theta)}\right\} .\label{eq:profit_emtrepreneur}
\end{equation}
\item Stop if $\max\left\{ \left|\mathbf{V}^{1}-\mathbf{V}^{0}\right|,\left|f^{1}-f^{0}\right|\right\} \leq\epsilon$.
Set $\mathbf{V}^{*}=\mathbf{V}^{1}$, and $f^{*}=f^{1}$. Otherwise,
set $\mathbf{V}^{0}=\mathbf{V}^{1}$, $f^{0}=f^{1}$ and repeat from
step (c). 
\end{enumerate}
\item Generate a large number of individuals, $N=100000$. For each agent
$j$ assign a vector of initial condition $\left(h_{t}^{j},z_{0}^{j},x_{0}^{j},a_{0}^{j},i_{w,0}^{j}\right)$,
where $x_{0}^{j}\sim\Omega_{x}$, $z_{0}^{j}\in\Omega_{z}$, $i_{w,0}^{j}\in\Omega_{w}$. 
\item Simulate the economy for $T$ periods, where $T$ is sufficiently
large. 
\item Calculate the following statistics from the simulated path $\left\{ h_{t}^{j},z_{t}^{j},x_{t}^{j},a_{t}^{j},i_{w,t}^{j},\mathcal{I}_{e,t}^{j},n_{t}^{j},k_{t}^{j}\right\} _{t=0}^{T}$.
\begin{align}
LS^{0} & =\frac{\sum_{j=1}^{N}\left(\iota_{i_{w}=1}\cdot z^{j}\right)}{N}\label{eq:agg_simulation_1}\\
KS^{0} & =\frac{\sum_{j=1}^{N}\left(\iota_{i_{w}=0}\cdot a^{j}+\iota_{i_{w}=1}\cdot a^{j}+\iota_{i_{w}=2}\cdot a^{j}-\iota_{i_{w}=0}\cdot k^{j}\right)}{\sum_{j=1}^{N}a^{j}}\label{eq:agg_simulation_2}\\
u_{t} & =\left(\int_{h}\iota_{i_{w}=0}+\int_{h}\iota_{i_{w}=2}\right)\cdot\int_{0}^{h_{u}^{*}}\\
v_{t} & =\int_{h}\iota_{i_{w}=1}\cdot\rho\\
p(\theta)_{t} & =\bar{m}\left(\frac{v_{t}}{u_{t}}\right)^{1-\mu}\\
G & =\int\tau wz
\end{align}
and $\tau_{y}^{1}$ that balances the government's budget. 
\item Stop and set $\left(r^{*},w^{*},\theta^{*}\right)=\left(r^{0},w^{0},\theta^{0}\right)$,
if $\max\left\{ LS^{0},KS^{0},\left|\theta_{n}^{1}-\theta_{n}^{0}\right|\right\} \leq\epsilon$.
Otherwise, update aggregate variables (restart from step 2): 
\begin{align*}
r^{0} & =\chi r^{0}+(1-\chi)r^{1}\left(KS^{0}\right)\\
w^{0} & =\chi w^{0}+(1-\chi)w^{1}\left(LS^{0}\right)\\
p(\theta)^{0} & =\chi p(\theta)^{0}+(1-\chi)p(\theta)^{1}\\
\tau_{y}^{0} & =\chi\tau_{y}^{0}+(1-\chi)\tau_{y}^{1}
\end{align*}
where $\chi\in(0,1)$ is the step for updating aggregate variables. 
\end{enumerate}

\section*{}

\section*{Appendix Search and Matching}

I am not sure how much of this is needed for coding purposes but this
is the standard search and matching model. The stock of employed workers
$n_{t}$ can be expressed as the current pool of workers who did not
separate from firms, plus new hires

\[
n_{t}=(1-\rho)n_{t-1}+m_{t-1}
\]
The matching function is Cobb-Douglass where

\[
m_{t}=\bar{m}v_{t}^{1-\mu}u_{t}^{\mu}
\]
$\bar{m}$ is a match efficiency parameter , $v$ is vacancies, $u$
is unemployment, and match elasticity is $0<\mu<1$.The hiring rate
can be defined for firms as total number of new hires to the number
of vacancies

\[
q(\theta_{t})=\frac{m_{t}}{v_{t}}=\bar{m}\theta_{t}^{-\mu}
\]
where $\theta_{t}=\frac{v_{t}}{u_{t}}$is aggregate labour market
tightness. Similarly we can define this from the perspective of the
worker in the job-finding rate.

\[
p(\theta_{t})=\frac{m_{t}}{u_{t}}=\bar{m}\theta_{t}^{1-\mu}
\]
The Job Creation Condition determines the amount of new vacancies
an employer will post. In our model, most workers are on rolling contracts,
so firms only need to post new vacancies after some workers separate
from these rolling contracts. Usually in search and matching models
the wage is determined ex-post of matching by Nash bargaining between
the employer and employee. However, in our span-of-control model the
wage has already been determined ex ante, as the wage is an endogenous
variable that determines occupational choice. The marginal benefit
that the new employee would bring to the firm is expressed as $A$
(which can be derived from the firms production function, $\gamma Xk^{\alpha}(zn)^{\gamma-1}$).
The $JCC\text{ can be then expressed as the costs of replacing out workers on the \ensuremath{LHS} and the expected benefit on the \ensuremath{RHS.}}$

\[
\frac{\kappa}{q(\theta_{t})}=A-w_{t}+\beta E_{t}\left[A-w_{t+1}+(1-\rho)\frac{\kappa}{q(\theta_{t+1})}\right]
\]
where $\kappa$ is the cost of posting a vacancy. %
\begin{comment}
This is different from the standard JCC as usually production starts
the period after matching so the JCC would be $\frac{\kappa}{q(\theta_{t})}=+\beta E_{t}\left[A-w_{t+1}+(1-\rho)\frac{\kappa}{q(\theta_{t+1})}\right].$
See alternative timing version to think about the problem this way,
I am not sure it makes much difference. 
\end{comment}


\subsection*{Matching Steady States}

To find the steady state of the employment equation $n_{t}=(1-\rho)n_{t-1}+m_{t-1}$,
we use the fact that $u=(1-n)$ and $m_{t}=\bar{m}v_{t}^{1-\mu}u_{t}^{\mu}.$
Substituting and getting rid of time subscript gives $(1-u)=(1-\rho)(1-u)+\bar{m}v^{1-\mu}u^{\mu}$
which simplifies to $\bar{m}v_{t}^{1-\mu}u_{t}^{\mu}=$$\rho(1-u)$.
Making $v\text{\,the subject of the foruma gives us the Beveridge curve.}$

\[
v=\left(\frac{\rho(1-u)}{\bar{mu}}\right)^{1-\mu}u
\]
The steady state for the JCC can be found by using $\frac{\kappa}{q(\theta_{t})}=A-w_{t}+\beta E_{t}\left[A-w_{t+1}+(1-\rho)\frac{\kappa}{q(\theta_{t+1})}\right]$
and $q(\theta_{t})=\bar{m}\theta_{t}^{-\mu}=\bar{m}\left(\frac{v_{t}}{u_{t}}\right)^{-\mu}$.
Substituting and getting rid of time subscripts yields 
\[
\frac{\kappa}{\bar{m}}\left(\frac{v}{u}\right)^{\mu}=A-w_{t}+\beta\left[A-w+(1-\rho)\frac{\kappa}{\bar{m}}\left(\frac{v}{u}\right)^{\mu}\right]
\]

Simplifying yields the steady state of the JCC

\[
1-\beta(1-\rho)\frac{\kappa}{\bar{m}}\left(\frac{v}{u}\right)^{\mu}=(1+\beta)\left[A-w\right]
\]

 \bibliographystyle{apalike}
\bibliography{../../../_shared/references/global_references}

\end{document}

